% Part: intuitionistic-logic
% Chapter: tableaux
% Section: introduction

\documentclass[../../../include/open-logic-section]{subfiles}

\begin{document}

\olfileid{int}{tab}{int}

\olsection{ভূমিকা}

!!^{tableau}s হলো !!{signed
  formula}s দিয়ে গঠিত বিশেষ
(নিচের দিকে শাখাবিভক্ত) বৃক্ষ। এর প্রতিটি সদস্য হলো
সত্যতার চিহ্ন ($\True$ অথবা $\False$) এবং !!a{sentence}-এর জোড়া:
\[
\sFmla{\True}{!A} \text{ অথবা } \sFmla{\False}{!A}.
\]
!!^a{tableau} শুরু হয় কয়েকটি \emph{অনুমান} দিয়ে। পরে প্রতিটি
!!{signed formula} পাওয়া যায় কোনো অনুমান-বিধি প্রয়োগ করে।
কিছু বিধি বৃক্ষের একটি প্রান্তে এক বা একাধিক !!{signed formula}s
যোগ করে; অন্যগুলি দুটি নতুন প্রান্ত যোগ করে, ফলে দুটি শাখা হয়।
একটি বিধি যে !!{signed formula}s দেয়, সেগুলির !!{formula}
বিধি-প্রয়োগের লক্ষ্য !!{signed formula}-র সূত্রের চেয়ে
কম জটিল। কোনো শাখায় $\sFmla{\True}{!A}$ এবং
$\sFmla{\False}{!A}$ উভয়ই থাকলে শাখাটিকে \emph{বন্ধ} বলি।
!!a{tableau}-র সব শাখা বন্ধ হলে পুরো !!{tableau} বন্ধ।
একটি বন্ধ !!{tableau} হলো এমন !!a{derivation} যা দেখায়,
!!{tableau} শুরুর !!{signed formula}s-এর সেটটি অসন্তোষযোগ্য।
এর সাহায্যে $\Proves$ সম্পর্ক সংজ্ঞায়িত করা যায়:
$\Gamma \Proves !A$ তখন এবং কেবল তখনই, যখন কোনো সসীম সেট
$\Gamma_0 = \{!B_1, \dots, !B_n\} \subseteq \Gamma$ আছে,
যার জন্য নিচের অনুমানগুলি দিয়ে একটি বন্ধ !!{tableau} হয়:
\[
\{\sFmla{\False}{!A}, \sFmla{\True}{!B_1}, \dots, \sFmla{\True}{!B_n}\}.
\]

স্বজ্ঞাবাদী যুক্তিবিদ্যার জন্য !!{signed formula}-র ধারণা
প্রসারিত এবং সংযোজকগুলির বিধি সংশোধিত করতে হবে।
% BN-SRC-409: এই অধ্যায় স্বজ্ঞাবাদী ট্যাবলোর; মোডাল ট্যাবলোর নয়।
চিহ্নের ($\True$ অথবা $\False$) পাশাপাশি স্বজ্ঞাবাদী
!!{tableau}s-এর !!{formula}s-তে \emph{উপসর্গ}~$\sigma$ থাকে।
উপসর্গগুলি ধনাত্মক পূর্ণসংখ্যার অশূন্য সসীম ক্রম; অর্থাৎ
$\sigma \in (\PosInt)^* \setminus \{\emptyseq\}$। এমন উপসর্গ
লিখতে উপসর্গ~$\sigma$-র চারপাশের $\tuple{\ }$ বাদ দিয়ে পৃথক !!{element}s-এর
মাঝে $,$-র বদলে~$.$ দিই। $\sigma$ উপসর্গ হলে
$\sigma.n$ হলো $\sigma \concat \tuple{n}$; যেমন,
$\sigma = 1.2.1$ হলে $\sigma.3$ হলো $1.2.1.3$।
উদাহরণস্বরূপ,
\[
\sFmla{\True}{!A \lif (!B \lif !C)}[1.2]
\]
একটি \emph{উপসর্গযুক্ত !!{signed formula}}
(সংক্ষেপে, \emph{উপসর্গযুক্ত !!{formula}})।

স্বজ্ঞাগতভাবে, উপসর্গটি এমন একটি মডেলের জগতের নাম দেয় যা
!!a{tableau}-র কোনো শাখার !!{formula}s সন্তুষ্ট করতে পারে।
$\sigma$ কোনো জগতের নাম হলে $\sigma.n$ সেই জগৎ থেকে
প্রবেশযোগ্য একটি জগতের নাম।

স্বজ্ঞাবাদী মডেলে প্রবেশযোগ্যতা-সম্পর্ক প্রতিবিম্বী ও পরিযায়ী।
উপসর্গের ভাষায়, $\sigma$ নিজে $\sigma$ থেকে প্রবেশযোগ্য;
$\sigma$-র যেকোনো সম্প্রসারণ, অর্থাৎ
$\sigma.n_1.\cdots.n_k$-রূপের উপসর্গও প্রবেশযোগ্য।
$\sigma$ নিজে ও তার সমস্ত সম্প্রসারণ বোঝাতে $\sigma.*$
লিখব। অন্যভাবে বললে, $\sigma.*$ উপসর্গগুলি ঠিক সেইগুলিই,
যেগুলি $\sigma$ থেকে প্রবেশযোগ্য।

\end{document}
